\documentclass[11pt]{article}

\usepackage[a4paper,margin=1in]{geometry}
\usepackage{amsmath,amssymb}
\usepackage{hyperref}

\hypersetup{
    colorlinks=true,
    linkcolor=blue,
    citecolor=blue,
    urlcolor=blue
}

\title{Steepest Descent Depends on the Metric}
\author{}
\date{}

\begin{document}

\maketitle

\section*{Steepest Among Which Directions?}

At a point $p$, the differential
\[
df_p:T_pM\to\mathbb R
\]
tells us the rate of change in every direction:
\[
X_p\mapsto df_p(X_p).
\]
If we want the direction of fastest decrease, we might try to minimize
\[
df_p(X_p)
\]
over all nonzero directions $X_p$.

But this has no answer unless we say how large a direction is. If $X_p$ gives
some decrease, then $100X_p$ gives one hundred times as much first-order
decrease:
\[
df_p(100X_p)=100\,df_p(X_p).
\]
So without a notion of unit length, ``steepest'' is meaningless.

\section*{The Metric Defines Unit Directions}

A Riemannian metric gives a norm
\[
\|X_p\|_g
=
\sqrt{g_p(X_p,X_p)}.
\]
Now we can ask a meaningful question:
\[
\text{among all }X_p\text{ with }\|X_p\|_g=1,
\text{ which minimizes }df_p(X_p)?
\]

Using
\[
df_p(X_p)
=
g_p(\operatorname{grad}_g f(p),X_p),
\]
the Cauchy--Schwarz inequality gives
\[
df_p(X_p)
\ge
-\|\operatorname{grad}_g f(p)\|_g\|X_p\|_g.
\]
For unit vectors,
\[
df_p(X_p)
\ge
-\|\operatorname{grad}_g f(p)\|_g.
\]
Equality occurs when
\[
X_p
=
-\frac{\operatorname{grad}_g f(p)}
{\|\operatorname{grad}_g f(p)\|_g}.
\]
Thus the negative gradient is the steepest descent direction among unit
directions measured by $g$.

\section*{No Metric, No Canonical Steepest Direction}

On a bare smooth manifold, there is no canonical norm on $T_pM$. Therefore
there is no canonical notion of unit direction, angle, or steepest descent.

You can still compute directional derivatives:
\[
df_p(X_p).
\]
But you cannot canonically say which direction is the steepest without choosing
some way to measure the size of directions.

\section*{Why Different Metrics Give Different Flows}

Changing the metric changes the unit sphere in each tangent space. Since
steepest descent means ``best decrease per unit length,'' changing the notion
of length changes the steepest direction.

This is why Euclidean gradient descent, natural gradient descent, Riemannian
gradient descent, and Wasserstein gradient flow can all have the same objective
but different dynamics: they use different metrics.

\end{document}
